% GENERATED FILE. Edit canonical lessons, then run npm run generate:books.
% canonical-source-hash: fnv1a64:f20bf6423fe9346e
% canonical-lessons: JA-C297-oni, JA-C297-omen, JA-C297-kowagaru, JA-C297-oidasu, JA-C297-mezurashii

\chapter{The Ogre at the Door}
\label{ch:ja-the-ogre-at-the-door}
\hlchaptermodality{297}

\begin{chapteropening}
\textbf{By the end of this chapter:} \emph{I can say how the ogre is driven out at the festival.}

Complete the last lesson of chapter 297: I can say how the ogre is driven out at the festival.
\end{chapteropening}

\section[\texorpdfstring{\ja{おに} (oni)}{おに (oni)}]{\ja{おに} (oni) — an ogre; a demon}

\label{lesson:JA-C297-oni}



\begin{quote}
Before the new one: say the Japanese for a custom, then the Japanese for to inherit.
\end{quote}

\subsection*{You'll want to know: \ja{おに}}
\textbf{\ja{おに}} — \emph{oni} — \textquotedblleft{}an ogre; a demon\textquotedblright{}.

\begin{quote}
\textbf{\ja{おには} \ja{こわい} \ja{かおを} \ja{しています。}} — \emph{oni wa kowai kao o shite imasu} — \textquotedblleft{}The ogre has a frightening face.\textquotedblright{}
\end{quote}

\begin{cousinweb}
\textbf{\ja{おに}} (\emph{oni}) is a native Japanese word (wago), inherited rather than borrowed from Chinese; its kanji was attached later as a kun'yomi (JMdict).
\end{cousinweb}

\begin{grammarlens}[title={What you've built}]
One more word for telling what happened.
\end{grammarlens}

\subsection*{Your turn}
\begin{itemize}
\raggedright
  \item \emph{Say it:} \emph{oni}
  \item \emph{Say it:} \emph{oni}, once more
  \item \emph{Say it:} \emph{oni wa kowai kao o shite imasu} — The ogre has a frightening face.
  \item \emph{Recall:} say the Japanese for a custom, then the Japanese for to inherit, then say \emph{oni} again
\end{itemize}

\begin{tcolorbox}[breakable,colback=teal!4,colframe=teal!35!black,title={Before you move on}]
What does \ja{おに} mean? (\textquotedblleft{}An ogre; a demon\textquotedblright{}.) Say it once more.
\end{tcolorbox}
\section[\texorpdfstring{\ja{おめん} (omen)}{おめん (omen)}]{\ja{おめん} (omen) — a mask}

\label{lesson:JA-C297-omen}



\begin{quote}
Before the new one: say the Japanese for to inherit, then the Japanese for an ogre.
\end{quote}

\subsection*{You'll want to know: \ja{おめん}}
\textbf{\ja{おめん}} — \emph{omen} — \textquotedblleft{}a mask\textquotedblright{}.

\begin{quote}
\textbf{\ja{あには} \ja{おにの} \ja{おめんを} \ja{かぶりました。}} — \emph{ani wa oni no omen o kaburimashita} — \textquotedblleft{}My brother put on an ogre mask.\textquotedblright{}
\end{quote}

\begin{cousinweb}
\textbf{\ja{おめん}} (\emph{omen}) is a mixed word, native and Sino-Japanese: \emph{o-}, \textquotedblleft{}the polite prefix\textquotedblright{}, is native Japanese, and \emph{men}, \textquotedblleft{}mask\textquotedblright{}, is a Chinese character read with its on'yomi, a reading borrowed from Chinese (Middle Chinese \emph{mjienH}) (JMdict, KANJIDIC2; Middle Chinese in Baxter's transcription).
\end{cousinweb}

\begin{grammarlens}[title={What you've built}]
One more word for telling what happened.
\end{grammarlens}

\subsection*{Your turn}
\begin{itemize}
\raggedright
  \item \emph{Say it:} \emph{omen}
  \item \emph{Say it:} \emph{omen}, once more
  \item \emph{Say it:} \emph{ani wa oni no omen o kaburimashita} — My brother put on an ogre mask.
  \item \emph{Recall:} say the Japanese for to inherit, then the Japanese for an ogre, then say \emph{omen} again
\end{itemize}

\begin{tcolorbox}[breakable,colback=teal!4,colframe=teal!35!black,title={Before you move on}]
What does \ja{おめん} mean? (\textquotedblleft{}A mask\textquotedblright{}.) Say it once more.
\end{tcolorbox}
\section[\texorpdfstring{\ja{こわがる} (kowagaru)}{こわがる (kowagaru)}]{\ja{こわがる} (kowagaru) — to be afraid; to be scared of}

\label{lesson:JA-C297-kowagaru}



\begin{quote}
Before the new one: say the Japanese for an ogre, then the Japanese for a mask.
\end{quote}

\subsection*{You'll want to know: \ja{こわがる}}
\textbf{\ja{こわがる}} — \emph{kowagaru} — \textquotedblleft{}to be afraid; to be scared of\textquotedblright{}.

\begin{quote}
\textbf{\ja{いもうとは} \ja{おにを} \ja{こわがっています。}} — \emph{imōto wa oni o kowagatte imasu} — \textquotedblleft{}My little sister is scared of the ogre.\textquotedblright{}
\end{quote}

Polite \textbf{\ja{こわがります}}, \emph{kowagarimasu}; a godan verb. Past (polite): \textbf{\ja{こどもが} \ja{こわがりました。}}, \emph{kodomo ga kowagarimashita}.

\begin{cousinweb}
\textbf{\ja{こわがる}} (\emph{kowagaru}) is a native Japanese word (wago), inherited rather than borrowed from Chinese, built from \emph{kowa}, \textquotedblleft{}fear\textquotedblright{} (from \emph{kowai}), and \emph{-garu}, \textquotedblleft{}to show signs of\textquotedblright{} (JMdict, KANJIDIC2).
\end{cousinweb}

\begin{grammarlens}[title={What you've built}]
One more word for telling what happened.
\end{grammarlens}

\subsection*{Your turn}
\begin{itemize}
\raggedright
  \item \emph{Say it:} \emph{kowagaru}
  \item \emph{Say it:} \emph{kowagaru}, once more
  \item \emph{Say it:} \emph{imōto wa oni o kowagatte imasu} — My little sister is scared of the ogre.
  \item \emph{Recall:} say the Japanese for an ogre, then the Japanese for a mask, then say \emph{kowagaru} again
\end{itemize}

\begin{tcolorbox}[breakable,colback=teal!4,colframe=teal!35!black,title={Before you move on}]
What does \ja{こわがる} mean? (\textquotedblleft{}To be afraid; to be scared of\textquotedblright{}.) Say it once more.
\end{tcolorbox}
\section[\texorpdfstring{\ja{おいだす} (oidasu)}{おいだす (oidasu)}]{\ja{おいだす} (oidasu) — to drive out; to chase away}

\label{lesson:JA-C297-oidasu}



\begin{quote}
Before the new one: say the Japanese for a mask, then the Japanese for to be afraid.
\end{quote}

\subsection*{You'll want to know: \ja{おいだす}}
\textbf{\ja{おいだす}} — \emph{oidasu} — \textquotedblleft{}to drive out; to chase away\textquotedblright{}.

\begin{quote}
\textbf{\ja{みんなで} \ja{おにを} \ja{おいだします。}} — \emph{minna de oni o oidashimasu} — \textquotedblleft{}We all drive the ogre out.\textquotedblright{}
\end{quote}

Polite \textbf{\ja{おいだします}}, \emph{oidashimasu}; a godan verb. Past (polite): \textbf{\ja{こどもが} \ja{おにを} \ja{おいだしました。}}, \emph{kodomo ga oni o oidashimashita}.

\begin{cousinweb}
\textbf{\ja{おいだす}} (\emph{oidasu}) is a native Japanese word (wago), inherited rather than borrowed from Chinese, built from \emph{oi}, \textquotedblleft{}chasing\textquotedblright{} (from \emph{ou}), and \emph{dasu}, \textquotedblleft{}to put out\textquotedblright{} (JMdict, KANJIDIC2).
\end{cousinweb}

\begin{grammarlens}[title={What you've built}]
One more word for telling what happened.
\end{grammarlens}

\subsection*{Your turn}
\begin{itemize}
\raggedright
  \item \emph{Say it:} \emph{oidasu}
  \item \emph{Say it:} \emph{oidasu}, once more
  \item \emph{Say it:} \emph{minna de oni o oidashimasu} — We all drive the ogre out.
  \item \emph{Recall:} say the Japanese for a mask, then the Japanese for to be afraid, then say \emph{oidasu} again
\end{itemize}

\begin{tcolorbox}[breakable,colback=teal!4,colframe=teal!35!black,title={Before you move on}]
What does \ja{おいだす} mean? (\textquotedblleft{}To drive out; to chase away\textquotedblright{}.) Say it once more.
\end{tcolorbox}
\section[\texorpdfstring{\ja{めずらしい} (mezurashii)}{めずらしい (mezurashii)}]{\ja{めずらしい} (mezurashii) — rare; unusual}

\label{lesson:JA-C297-mezurashii}



\begin{quote}
Before the new one: say the Japanese for to be afraid, then the Japanese for to drive out.
\end{quote}

\subsection*{You'll want to know: \ja{めずらしい}}
\textbf{\ja{めずらしい}} — \emph{mezurashii} — \textquotedblleft{}rare; unusual\textquotedblright{}.

\begin{quote}
\textbf{\ja{これは} \ja{めずらしい} \ja{まつりです。}} — \emph{kore wa mezurashii matsuri desu} — \textquotedblleft{}This is an unusual festival.\textquotedblright{}
\end{quote}

An i-adjective.

\begin{cousinweb}
\textbf{\ja{めずらしい}} (\emph{mezurashii}) is a native Japanese word (wago), inherited rather than borrowed from Chinese; its kanji was attached later as a kun'yomi (JMdict).
\end{cousinweb}

\begin{grammarlens}[title={What you've built}]
One more word for telling what happened.
\end{grammarlens}

\subsection*{Your turn}
\begin{itemize}
\raggedright
  \item \emph{Say it:} \emph{mezurashii}
  \item \emph{Say it:} \emph{mezurashii}, once more
  \item \emph{Say it:} \emph{kore wa mezurashii matsuri desu} — This is an unusual festival.
  \item \emph{Recall:} say the Japanese for to be afraid, then the Japanese for to drive out, then say \emph{mezurashii} again
\end{itemize}

\begin{tcolorbox}[breakable,colback=teal!4,colframe=teal!35!black,title={Before you move on}]
What does \ja{めずらしい} mean? (\textquotedblleft{}Rare; unusual\textquotedblright{}.) Say it once more.
\end{tcolorbox}
